\documentclass[11pt,a4paper]{article}
\usepackage[italian]{babel}
\usepackage[margin=2.5cm]{geometry}
\usepackage{parskip}
\usepackage{amsmath, amssymb, amsthm}
\usepackage{booktabs}
\usepackage{array}
\usepackage{xcolor}
\usepackage{needspace}
\usepackage{enumitem}
\usepackage{listings}
\usepackage{tikz}
\usepackage[most]{tcolorbox}
\usepackage[hidelinks]{hyperref}
\usetikzlibrary{decorations.pathreplacing, shapes.geometric, arrows.meta, positioning}

% Niente vedove/orfane: i paragrafi non lasciano righe isolate a fine/inizio pagina
\widowpenalty=10000
\clubpenalty=10000

% Elenchi più compatti
\setlist{itemsep=2pt, parsep=0pt, topsep=4pt, partopsep=0pt}

% --- Riquadri con bordo nero, senza numero (si spezzano tra due pagine) ---
% Uso: \begin{definizione} ... \end{definizione}
%      \begin{teorema}[Nome] ... \end{teorema}  (il nome è facoltativo)
\tcbset{nota/.style={enhanced jigsaw, breakable, colback=white,
  colframe=black, boxrule=0.5pt, sharp corners}}
\NewTColorBox{definizione}{}{nota,
  before upper={\textbf{Definizione.}\ }}
\NewTColorBox{teorema}{o}{nota,
  before upper={\textbf{Teorema\IfValueT{#1}{ (#1)}.}\ }}
\NewTColorBox{proposizione}{o}{nota,
  before upper={\textbf{Proposizione\IfValueT{#1}{ (#1)}.}\ }}
\NewTColorBox{proprieta}{o}{nota,
  before upper={\textbf{Proprietà\IfValueT{#1}{ (#1)}.}\ }}
\NewTColorBox{assioma}{o}{nota, boxrule=1pt,
  before upper={\textbf{Assioma\IfValueT{#1}{ (#1)}.}\ }}

% --- Ambienti semplici (senza riquadro) ---
% Il titolo sta in cima al blocco, anche se il contenuto inizia con un riquadro o
% con codice; e non resta mai da solo in fondo a una pagina.
% Uso: \begin{esempio}[Titolo facoltativo] ... \end{esempio}
\newcommand{\newrunin}[2]{%
  \NewDocumentEnvironment{#1}{o}{%
    \par\Needspace{4\baselineskip}\addvspace{\medskipamount}\noindent
    \textbf{#2}\IfValueT{##1}{ (##1)}\textbf{.}\ \ignorespaces}%
  {\par\addvspace{\medskipamount}}}
\newrunin{esempio}{Esempio}
\newrunin{esercizio}{Esercizio}
\NewTColorBox{osservazione}{}{nota,
  before upper={\textbf{Osservazione.}\ }}

% V e F nelle tavole di verità
\newcommand{\V}{\textsf{V}}
\newcommand{\F}{\textsf{F}}

% --- Codice C ---
% Uso: \begin{codice} ... \end{codice}      codice C numerato
%      \begin{comando} ... \end{comando}    comandi da terminale
%      \begin{risultato} ... \end{risultato} output di un programma
%      \lstinline|printf("ciao");|           codice dentro al testo
\lstdefinestyle{c}{
  language=C,
  basicstyle=\ttfamily\small,
  keywordstyle=\color{blue}\bfseries,
  commentstyle=\color{gray}\itshape,
  stringstyle=\color{red!70!black},
  numbers=left,
  numberstyle=\tiny\color{gray},
  numbersep=8pt,
  columns=flexible,
  keepspaces=true,
  breaklines=true,
  showstringspaces=false,
  tabsize=4,
  morekeywords={include, define},
  % lettere accentate nei commenti
  literate={à}{{\`a}}1 {è}{{\`e}}1 {é}{{\'e}}1 {ì}{{\`i}}1
           {ò}{{\`o}}1 {ù}{{\`u}}1 {È}{{\`E}}1 {À}{{\`A}}1
}
\lstset{style=c}
\newtcblisting{codice}{listing only, nota, left=6mm,
  listing options={style=c}}
\newtcblisting{comando}{listing only, nota,
  listing options={basicstyle=\ttfamily\small, numbers=none, language={},
    columns=flexible, keepspaces=true, breaklines=true}}
\newtcblisting{risultato}{listing only, enhanced jigsaw, breakable,
  colback=black!5, colframe=black, boxrule=0.5pt, sharp corners,
  listing options={basicstyle=\ttfamily\small, numbers=none, language={},
    columns=flexible, keepspaces=true, breaklines=true}}

% --- Diagrammi di flusso (simboli standard) ---
\tikzset{
  terminale/.style = {ellipse, draw, minimum width=2.5cm, minimum height=0.9cm},
  io/.style        = {trapezium, trapezium left angle=70, trapezium right angle=110,
                      draw, minimum width=2.5cm, minimum height=0.9cm},
  processo/.style  = {rectangle, draw, minimum width=2.5cm, minimum height=0.9cm},
  decisione/.style = {diamond, draw, aspect=2, minimum width=2.5cm},
  freccia/.style   = {thick, -{Stealth}}
}

\title{Bistabili, latch, flip-flop e sintesi di circuiti sequenziali}
\author{Marco Cerbella}
\date{Lezione 8}

\begin{document}
\maketitle

\section{Elementi bistabili}

Il blocco costitutivo fondamentale di un registro di memoria è un elemento \emph{bistabile}: un elemento circuitale con due stati stabili, in grado dunque di memorizzare un bit di informazione. Un elemento con $N$ stati stabili conserva $\log_2 N$ bit.

\begin{definizione}[Stato stabile]
Dato un circuito (sequenziale), si dice \emph{stato stabile} del circuito un'attribuzione di valori di verità a tutti i fili del circuito in accordo con il comportamento delle porte logiche.
\end{definizione}

\subsection{Due negatori collegati a croce}

Il bistabile più semplice è formato da una coppia di negatori $I_1$, $I_2$ collegati a croce: l'ingresso di $I_1$ è l'uscita di $I_2$ e viceversa.

\begin{center}
\begin{tikzpicture}[>=Stealth, every node/.style={font=\small}]
  \node[draw, regular polygon, regular polygon sides=3, shape border rotate=-90, minimum size=1cm, inner sep=0] (i1) at (0,1) {};
  \node[draw, regular polygon, regular polygon sides=3, shape border rotate=90, minimum size=1cm, inner sep=0] (i2) at (0,-0.4) {};
  \node at (-0.15,1) {$I_1$};
  \node at (0.15,-0.4) {$I_2$};
  \draw (i1.east) -- ++(1.2,0) coordinate (q) node[right] {$Q$};
  \draw (i2.west) -- ++(-1.2,0) coordinate (qn) node[left] {$\overline{Q}$};
  \draw (i1.west) -- ++(-0.5,0) |- (i2.west |- qn) ;
  \draw (i2.east) -- ++(0.5,0) |- (i1.east |- q);
\end{tikzpicture}
\end{center}

L'analisi di un elemento sequenziale differisce da quella di un circuito combinatorio perché il primo è \emph{ciclico}: le uscite $Q$ e $\overline{Q}$ dipendono l'una dall'altra. Si studiano i due casi $Q = 0$ e $Q = 1$.

\begin{itemize}
  \item \textbf{Caso $Q = 0$.} $I_2$ riceve in ingresso $0$ e produce $1$ su $\overline{Q}$; $I_1$ riceve $1$ da $\overline{Q}$ e produce $0$ su $Q$. Lo stato è stabile perché viene confermata la premessa iniziale ($Q = 0$).
  \item \textbf{Caso $Q = 1$.} $I_2$ riceve $1$ e produce $0$ su $\overline{Q}$; $I_1$ riceve $0$ e produce $1$ su $Q$. Anche questo stato è stabile.
\end{itemize}

\begin{osservazione}
Il circuito ha in realtà anche un terzo possibile stato, con entrambe le uscite circa a metà tra $0$ e $1$: è detto stato \emph{metastabile}. Inoltre, la variabile di stato è una sola: $\overline{Q}$ non contiene informazione aggiuntiva, essendo il complemento di $Q$.
\end{osservazione}

\paragraph{Limiti.} Due negatori collegati a croce memorizzano un bit, e un insieme di coppie di negatori potrebbe codificare le variabili di stato di una macchina a stati finiti; tuttavia l'utente non ha a disposizione \emph{alcun ingresso} per controllare tali variabili. I latch e i flip-flop che seguono sono provvisti di ingressi per controllare il valore delle variabili di stato, e sono quindi più utili per codificare gli stati di una FSM.

\section{Latch}

\subsection{Latch SR}

Il latch SR si realizza con due porte NOR $N_1$, $N_2$ collegate a croce. Ha due ingressi $S$ (\emph{Set}) e $R$ (\emph{Reset}) e due uscite $Q$ e $\overline{Q}$.

\begin{center}
\begin{tikzpicture}[>=Stealth, every node/.style={font=\small}]
  \node[draw, minimum width=1.6cm, minimum height=0.8cm] (n1) at (0,1) {NOR $N_1$};
  \node[draw, minimum width=1.6cm, minimum height=0.8cm] (n2) at (0,-1) {NOR $N_2$};
  \draw (-0.8,1.2) -- (-1.6,1.2) node[left] {$R$};
  \draw (-0.8,-1.2) -- (-1.6,-1.2) node[left] {$S$};
  \draw (0.8,1) -- (1.8,1) node[right] {$Q$};
  \draw (0.8,-1) -- (1.8,-1) node[right] {$\overline{Q}$};
  \fill (1.2,1) circle (1.5pt);
  \fill (1.2,-1) circle (1.5pt);
  \draw (1.2,1) -- (1.2,0) -- (-1.2,0) -- (-1.2,-0.8) -- (-0.8,-0.8);
  \draw (1.0,-1) -- (1.0,0.3) -- (-1.0,0.3) -- (-1.0,0.8) -- (-0.8,0.8);
\end{tikzpicture}
\end{center}

\paragraph{Analisi.}
\begin{itemize}
  \item \textbf{$S = 0$, $R = 0$.} Nessun ingresso è attivo: l'uscita $Q$ ricorda il suo valore precedente $Q_{\mathrm{prec}}$ (e $\overline{Q}$ il suo complemento). Vi sono due stati stabili, a seconda di $Q_{\mathrm{prec}}$.
  \item \textbf{$S = 1$, $R = 0$.} $N_2$ riceve almeno un valore vero ($S = 1$) e produce $0$ su $\overline{Q}$; a questo punto $N_1$ ha entrambi gli ingressi a $0$ e produce $Q = 1$. Quando viene attivato $S$, $Q$ viene \emph{settato} a $1$.
  \item \textbf{$S = 0$, $R = 1$.} $N_1$ riceve almeno un valore vero ($R = 1$) e produce $0$ su $Q$; $N_2$ ha entrambi gli ingressi a $0$ e produce $\overline{Q} = 1$. Quando viene attivato $R$, $Q$ viene \emph{resettato} a $0$.
  \item \textbf{$S = 1$, $R = 1$.} $N_1$ e $N_2$ vedono entrambe almeno un ingresso a $1$ e producono entrambe $0$. Non ha senso: significa settare e resettare contemporaneamente, e $Q$ e $\overline{Q}$ non sono l'una l'opposto dell'altra (inconsistenti).
\end{itemize}

\begin{center}
\begin{tabular}{cc|cc}
\toprule
$S$ & $R$ & $Q$ & $\overline{Q}$ \\
\midrule
0 & 0 & $Q_{\mathrm{prec}}$ & $\overline{Q}_{\mathrm{prec}}$ \\
0 & 1 & 0 & 1 \\
1 & 0 & 1 & 0 \\
1 & 1 & 0 & 0 \\
\bottomrule
\end{tabular}
\end{center}

\paragraph{Limiti.} Il latch SR è scomodo per due motivi:
\begin{itemize}
  \item si comporta in modo inconsistente quando $S$ e $R$ sono attivati contemporaneamente (entrambe le uscite valgono $0$);
  \item gli ingressi $S$ e $R$ combinano gli aspetti del \emph{come} e del \emph{quando}: attivare uno di essi determina non solo come debba diventare lo stato $Q$, ma anche quando debba cambiare.
\end{itemize}
Questi limiti sono superati dal latch D.

\subsection{Latch D}

Il latch D ha due ingressi:
\begin{itemize}
  \item un ingresso dati $D$, che controlla il \emph{prossimo stato};
  \item un ingresso di clock $CLK$, che controlla il \emph{momento} del cambio di stato.
\end{itemize}
Si ottiene da un latch SR ponendo $S = D \cdot CLK$ e $R = \overline{D} \cdot CLK$.

\begin{center}
\begin{tikzpicture}[>=Stealth, every node/.style={font=\small}]
  \node[draw, minimum width=1.0cm, minimum height=0.7cm] (a1) at (2,0.9) {AND};
  \node[draw, minimum width=1.0cm, minimum height=0.7cm] (a2) at (2,-0.9) {AND};
  \node[draw, minimum width=0.8cm, minimum height=0.5cm] (inv) at (0.4,0.7) {NOT};
  \node[draw, minimum width=1.6cm, minimum height=2.2cm] (sr) at (4.4,0) {latch SR};
  \draw[->] (a1.east) -- (3.6,0.9) node[pos=0.8, above, font=\footnotesize] {$R$};
  \draw[->] (a2.east) -- (3.6,-0.9) node[pos=0.8, below, font=\footnotesize] {$S$};
  \draw (5.2,0.9) -- (6.0,0.9) node[right] {$Q$};
  \draw (5.2,-0.9) -- (6.0,-0.9) node[right] {$\overline{Q}$};
  \draw (-1.2,-1.1) node[left] {$D$} -- (1.5,-1.1);
  \fill (-0.4,-1.1) circle (1.5pt);
  \draw (-0.4,-1.1) -- (-0.4,0.7) -- (inv.west);
  \draw (inv.east) -- (1.5,0.7);
  \draw (-1.2,1.9) node[left] {$CLK$} -- (0.9,1.9);
  \fill (0.9,1.9) circle (1.5pt);
  \draw (0.9,1.9) -- (0.9,-0.7) -- (1.5,-0.7);
  \draw (0.9,1.1) -- (1.5,1.1);
  \fill (0.9,1.1) circle (1.5pt);
\end{tikzpicture}
\end{center}

\paragraph{Analisi.}
\begin{center}
\begin{tabular}{cc|ccc|cc}
\toprule
$CLK$ & $D$ & $\overline{D}$ & $S$ & $R$ & $Q$ & $\overline{Q}$ \\
\midrule
0 & $\times$ & $\overline{\times}$ & 0 & 0 & $Q_{\mathrm{prec}}$ & $\overline{Q}_{\mathrm{prec}}$ \\
1 & 0 & 1 & 0 & 1 & 0 & 1 \\
1 & 1 & 0 & 1 & 0 & 1 & 0 \\
\bottomrule
\end{tabular}
\end{center}
Quando $CLK = 0$ il latch ricorda lo stato precedente; quando $CLK = 1$ copia $D$ su $Q$. Il caso $S = R = 1$ non può mai presentarsi, perché $S$ e $R$ sono sempre uno il complemento dell'altro quando $CLK = 1$.

\paragraph{Limiti.} Il latch D aggiorna \emph{continuamente} il suo stato finché $CLK = 1$, mentre sarebbe utile aggiornare lo stato solo in un momento preciso. Questo limite è superato dai flip-flop.

\section{Flip-flop D}

Un flip-flop D si realizza con due latch D in cascata, controllati da due segnali di clock complementari: il primo latch $L_1$ è detto \emph{master}, il secondo $L_2$ \emph{slave}; il nodo $N_1$ collega i due latch.

\begin{center}
\begin{tikzpicture}[>=Stealth, every node/.style={font=\small}]
  \node[draw, minimum width=1.4cm, minimum height=1.2cm] (l1) at (0,0) {$L_1$ (master)};
  \node[draw, minimum width=1.4cm, minimum height=1.2cm] (l2) at (3.2,0) {$L_2$ (slave)};
  \draw[->] (-1.6,0) node[left] {$D$} -- (l1.west) node[pos=0.5, above, font=\footnotesize] {};
  \draw[->] (l1.east) -- (l2.west) node[midway, above, font=\footnotesize] {$N_1$};
  \draw[->] (l2.east) -- ++(1.2,0) node[right] {$Q$};
  \draw[->] (-1.6,1.4) node[left] {$CLK$} -- (l1.north |- 0,1.4) -- (l1.north);
  \draw (0.0,1.4) -- (1.0,1.4) node[draw, circle, inner sep=1.5pt, fill=white] (c) {} (c) -- (3.2,1.4) edge[->] (l2.north);
\end{tikzpicture}
\end{center}

\paragraph{Analisi.}
\begin{itemize}
  \item Quando $CLK = 0$, il master è trasparente e qualsiasi valore di $D$ viene portato in $N_1$; lo slave è opaco e mantiene l'uscita.
  \item Quando $CLK$ diventa $1$, il master si blocca e il valore di $N_1$ viene trasmesso sull'uscita dell'elemento ($Q$ in $L_2$).
\end{itemize}
In altre parole, un flip-flop D \emph{copia $D$ su $Q$ al fronte di salita del clock} e ricorda il suo stato in tutti gli altri casi.

Il simbolo circuitale del flip-flop D è un rettangolo con ingressi $D$ e $CLK$ (quest'ultimo contrassegnato da un triangolo, che indica la sensibilità al fronte) e uscite $Q$ e $\overline{Q}$; esiste una versione semplificata con la sola uscita $Q$.

\section{Registro di memoria}

\begin{definizione}[Registro a $N$ bit]
Un \emph{registro} a $N$ bit è un banco di $N$ flip-flop che condividono un ingresso di clock $CLK$ in comune, in modo che tutti i bit vengano aggiornati allo stesso tempo.
\end{definizione}

\section{Sintesi di un circuito sequenziale: esempio}

\begin{esempio}[Semaforo di un incrocio]
\textbf{Specifiche.}
\begin{itemize}
  \item Incrocio tra due strade, nord-sud (NS) ed est-ovest (EO), controllate da un semaforo.
  \item Il semaforo può commutare ogni $30$ secondi; per semplicità si considerano solo rosso e verde.
  \item Un sensore legge, per ogni direttrice, se esiste almeno un'auto in attesa oppure un'auto che si accinga ad attraversare (condizioni trattate allo stesso modo).
  \item Il semaforo deve cambiare colore, da rosso a verde, quando esiste un'auto in attesa sulla sua direttrice.
  \item Se ci sono auto in attesa su entrambe le direttrici il semaforo deve cambiare colore (al termine del tempo di commutazione).
\end{itemize}

\textbf{Codifica.} Il primo passo è modellare il circuito con una macchina a stati finiti.
\begin{itemize}
  \item Gli ingressi sono codificati con due variabili booleane $i_{EO}$, $i_{NS}$.
  \item Serve $1$ bit per codificare lo stato (due stati possibili): serve dunque un solo bistabile, un flip-flop D $F$ con uscite $Q_F$, $\overline{Q}_F$. Si pone $Q_F = 1$ per lo stato VerdeNS e $Q_F = 0$ per VerdeEO.
  \item Serve $1$ bit per l'uscita: $u = 1$ codifica LuceVerdeNS, $u = 0$ codifica LuceVerdeEO.
\end{itemize}

\textbf{Tabella di transizione.} $D_F$ è lo stato futuro.
\begin{center}
\begin{tabular}{c|cc|cc}
\toprule
$Q_F$ & $i_{EO}$ & $i_{NS}$ & $D_F$ & $u$ \\
\midrule
1 & 0 & 0 & 1 & 1 \\
1 & 0 & 1 & 1 & 1 \\
1 & 1 & 0 & 0 & 1 \\
1 & 1 & 1 & 0 & 1 \\
0 & 0 & 0 & 0 & 0 \\
0 & 0 & 1 & 1 & 0 \\
0 & 1 & 0 & 0 & 0 \\
0 & 1 & 1 & 1 & 0 \\
\bottomrule
\end{tabular}
\end{center}

\textbf{Mappe di Karnaugh.} Dalla tabella si ottengono
\[
D_F = \overline{Q}_F\, i_{NS} + Q_F\, \overline{i_{EO}}, \qquad u = Q_F .
\]
Infatti, per $Q_F = 0$ lo stato futuro vale $1$ solo se $i_{NS} = 1$ (cioè $D_F = i_{NS}$), mentre per $Q_F = 1$ vale $1$ solo se $i_{EO} = 0$ (cioè $D_F = \overline{i_{EO}}$). L'uscita coincide con lo stato.

\textbf{Circuito.} Il circuito si sintetizza con due porte AND, una porta OR e un negatore: il flip-flop D $F$ ha in ingresso $D_F = \overline{Q}_F\, i_{NS} + Q_F\, \overline{i_{EO}}$ e l'uscita $u$ è $Q_F$.

\begin{center}
\begin{tikzpicture}[>=Stealth, every node/.style={font=\small}]
  \node[draw, minimum width=1.0cm, minimum height=0.7cm] (a1) at (2,1) {AND};
  \node[draw, minimum width=1.0cm, minimum height=0.7cm] (a2) at (2,-1) {AND};
  \node[draw, minimum width=0.9cm, minimum height=0.8cm] (o) at (4,0) {OR};
  \node[draw, minimum width=1.2cm, minimum height=1.8cm] (ff) at (6,0) {$F$};
  \node[draw, minimum width=0.8cm, minimum height=0.5cm] (n) at (0.2,0.8) {NOT};
  \draw (-1.2,0.8) node[left] {$i_{EO}$} -- (n.west);
  \draw (n.east) -- (1.5,0.8);
  \draw (-1.2,-0.8) node[left] {$i_{NS}$} -- (1.5,-0.8);
  \draw (2.5,1) -- (3.0,1) -- (3.0,0.2) -- (3.55,0.2);
  \draw (2.5,-1) -- (3.0,-1) -- (3.0,-0.2) -- (3.55,-0.2);
  \draw[->] (4.45,0) -- (5.4,0) node[midway, above, font=\footnotesize] {$D_F$};
  \draw (6.6,0.5) -- (7.8,0.5) node[right] {$u = Q_F$};
  \draw (6.6,-0.5) -- (7.0,-0.5);
  \fill (7.0,0.5) circle (1.5pt);
  \draw (7.0,0.5) -- (7.0,2.2) -- (1.0,2.2) -- (1.0,1.2) -- (1.5,1.2);
  \draw (7.0,-0.5) -- (7.0,-2.2) -- (1.0,-2.2) -- (1.0,-1.2) -- (1.5,-1.2);
  \node[font=\footnotesize] at (5.6,0.5) {};
  \node[font=\footnotesize, above] at (4,2.2) {$Q_F$};
  \node[font=\footnotesize, below] at (4,-2.2) {$\overline{Q}_F$};
\end{tikzpicture}
\end{center}
\end{esempio}

\begin{osservazione}
In sintesi, il procedimento per realizzare un circuito sequenziale è: modellare il problema con una FSM; codificare ingressi, stati e uscite con variabili booleane; scrivere la tabella di transizione (stato futuro e uscita in funzione di stato presente e ingressi); semplificare con le mappe di Karnaugh; realizzare la rete combinatoria e collegarla a un registro di flip-flop D che memorizza lo stato.
\end{osservazione}

\paragraph{Riferimenti.} Capitolo 3 del libro di supporto (S.\ Harris et al., \emph{Sistemi Digitali e Architettura dei Calcolatori}), sezioni 3.1, 3.2, 3.2.1, 3.2.2, 3.2.3 e 3.2.4.

\end{document}
